% Generated by maint/texbuild/theory_tex.py from theory/workbooks/cell_tracking/walk_back.md.
% Edits here are overwritten. Edit the markdown, then run the script again.
\chapter{\texorpdfstring{Walk back}{Walk back}}
\label{chap:walk-back}

\textbf{Purpose:} The rule every tracker stage is held to: it can be walked back to what it came from.

\textbf{Scope:} every stage of the driver (\hyperref[chap:cell-tracking-table]{cell\_\allowbreak{}tracking\_\allowbreak{}table.md}, its driver section).

Ruled by Doug, 25 September: “The whole point of the tracker is to track cells. If you can't walk back something you did you built it wrong.”

\begin{itemize}
\item Every stage keeps what it came from: a point keeps its frame, voxel and exact level; a transform keeps its exact inverse; a link keeps its two points and its evidence; a choice in the sort keeps the options it passed over, with their weights.
\item Nothing is deleted. What the sort does not keep is marked, and the mark can be walked back.
\item A map that sends two readings to one (a clip, a rounding, a bin) cannot be walked back and is not built.
\item What prompted it: OrganoidTracker's 1\%/99\% intensity clip was proposed for S0. It was replaced by a per-sample affine footing held as a pair, plus the set's cumulative count as the contrast; both are invertible.
\item It is written into the driver section of \hyperref[chap:cell-tracking-table]{cell\_\allowbreak{}tracking\_\allowbreak{}table.md}.
\end{itemize}
